\documentclass[11pt]{article}
\usepackage{amsmath,amssymb,amsthm}
\usepackage[margin=1in]{geometry}
\newtheorem{theorem}{Theorem}
\newtheorem{lemma}[theorem]{Lemma}
\title{Problem 364: Comfortable Distance}
\author{Project Euler Solutions}
\date{}
\begin{document}
\maketitle

\section{Problem Statement}

A row of $n$ seats is filled one person at a time. Each person chooses the seat maximizing their minimum distance to any occupied seat (leftmost in case of ties; the first person takes an end seat). Let $f(n)$ be the number of distinct valid filling orderings. Compute $f(n) \bmod 10^8$.

\section{Mathematical Foundation}

\begin{theorem}[Gap Splitting]
When a person sits in a contiguous gap of $L$ empty seats, they occupy position $m = \lceil L/2 \rceil$ from the left boundary, producing independent sub-gaps of sizes $m-1$ and $L-m$.
\end{theorem}

\begin{proof}
The greedy rule maximizes the minimum distance to occupied neighbors. In a gap of $L$ seats, the midpoint $\lceil L/2 \rceil$ achieves this maximum. The resulting two sub-gaps share no seats and are filled independently.
\end{proof}

\begin{lemma}[Interleaving Count]
If $k$ independent sub-gaps require $n_1, \ldots, n_k$ people, the number of valid filling orderings across all gaps is
\[
\binom{n_1 + \cdots + n_k}{n_1,\, n_2,\, \ldots,\, n_k}.
\]
\end{lemma}

\begin{proof}
Within each sub-gap, the filling order is determined by the greedy rule. The only choice is which sub-gap to fill at each time step when multiple gaps tie. This is equivalent to interleaving $k$ fixed sequences, giving the multinomial coefficient.
\end{proof}

\begin{theorem}[Recursive Formula]
Let $C(L)$ be the number of valid orderings for a gap of size~$L$. Then $C(0) = C(1) = 1$ and for $L \ge 2$:
\[
C(L) = \binom{L-1}{\lceil L/2 \rceil - 1} \cdot C\!\left(\lceil L/2 \rceil - 1\right) \cdot C\!\left(L - \lceil L/2 \rceil\right).
\]
\end{theorem}

\begin{proof}
The first placement in a gap of size~$L$ goes to position $m = \lceil L/2\rceil$, leaving sub-gaps of sizes $a = m-1$ and $b = L - m$ with $a + b = L - 1$ remaining people. The interleaving count is $\binom{L-1}{a}$. Each sub-gap contributes $C(a)$ and $C(b)$ internal orderings.
\end{proof}

\begin{lemma}[First Person]
The first person sits at seat~1, creating one gap of size $n-1$. Hence $f(n) = C(n-1)$.
\end{lemma}

\section{Algorithm}

\begin{verbatim}
function solve(n):
    M = 10^8
    precompute fact[0..n], inv_fact[0..n] mod M
    memo = {}
    function C(L):
        if L <= 1: return 1
        if L in memo: return memo[L]
        m = ceil(L / 2)
        result = binom(L-1, m-1) * C(m-1) * C(L-m) mod M
        memo[L] = result
        return result
    return C(n - 1)
\end{verbatim}

\section{Complexity Analysis}

\begin{itemize}
\item \textbf{Time:} $O(n)$ for factorials; $O(\log n)$ for the recursion (gap sizes halve).
\item \textbf{Space:} $O(n)$ for factorials; $O(\log n)$ for memoization.
\end{itemize}

\section{Answer}

\[
\boxed{44855254}
\]

\end{document}
